\documentclass{article}

\usepackage{amsmath,amssymb}
\usepackage{xurl}
\usepackage[colorlinks=true,linkcolor=blue,citecolor=blue,urlcolor=blue]{hyperref}
\setlength{\emergencystretch}{2em}

% Shared commands for notes in docs/paper-gaps/.

\DeclareUnicodeCharacter{2080}{\ifmmode{}_{0}\else\textsubscript{0}\fi}
\DeclareUnicodeCharacter{2081}{\ifmmode{}_{1}\else\textsubscript{1}\fi}
\DeclareUnicodeCharacter{2082}{\ifmmode{}_{2}\else\textsubscript{2}\fi}
\DeclareUnicodeCharacter{2099}{\ifmmode{}_{n}\else\textsubscript{n}\fi}

\newcommand{\Lean}{\textsc{Lean}}
\ExplSyntaxOn
\NewDocumentCommand{\leanid}{m}{
  \ifmmode
    \textsf{\detokenize{#1}}
  \else
    \group_begin:
    \urlstyle{sf}
    \tl_set:Nn \l_tmpa_tl {#1}
    \tl_replace_all:Nnn \l_tmpa_tl {\_} {_}
    \exp_args:NV \path \l_tmpa_tl
    \group_end:
  \fi
}
\ExplSyntaxOff
\providecommand{\tr}{\operatorname{tr}}
\newcommand{\tnleanIssueBase}{https://github.com/LionSR/TNLean/issues/}
\newcommand{\tnleanPrBase}{https://github.com/LionSR/TNLean/pull/}
\newcommand{\tnleanDocBase}{https://sirui-lu.com/TNLean/docs/}
\newcommand{\ghissue}[1]{\href{\tnleanIssueBase#1}{GitHub issue \##1}}
\newcommand{\ghpr}[1]{\href{\tnleanPrBase#1}{PR \##1}}
\newcommand{\doclink}[1]{\href{\tnleanDocBase#1.html}{\path{#1}}}

% Dirac notation and the identity operator, as in blueprint/src/macros/common.tex.
% \providecommand keeps note-local definitions authoritative.
\usepackage{dsfont}
\providecommand{\ket}[1]{|#1\rangle}
\providecommand{\bra}[1]{\langle #1|}
\providecommand{\braket}[2]{\langle #1 | #2 \rangle}
\providecommand{\ketbra}[2]{|#1\rangle\!\langle #2|}
\providecommand{\Id}{\mathds{1}}
\providecommand{\one}{\mathds{1}}
\providecommand{\C}{\mathbb{C}}
\providecommand{\MN}[1]{M_{#1}(\C)}
\providecommand{\MD}{M_D(\C)}

% External referencing for paper sources.  The build helper writes sanitized
% auxiliary files under build/paper-aux/ and excludes duplicated source labels.
\usepackage{xr}

% Import a generated paper auxiliary file with a caller-supplied label prefix.
% The candidate paths support builds from docs/paper-gaps/, the repository root,
% and build/paper-aux/ itself.
\newcommand{\importpaperaux}[2]{%
  \IfFileExists{../../build/paper-aux/#2.aux}{%
    \externaldocument[#1]{../../build/paper-aux/#2}%
  }{%
    \IfFileExists{build/paper-aux/#2.aux}{%
      \externaldocument[#1]{build/paper-aux/#2}%
    }{%
      \IfFileExists{#2.aux}{%
        \externaldocument[#1]{#2}%
      }{}%
    }%
  }%
}

% General external reference macros
\newcommand{\paperref}[2]{\ref{#1-#2}}
\newcommand{\papereqref}[2]{(\ref{#1-#2})}

% CPSV16 source (1606.00608 / MPDO-22-12-17-2.tex) external document and shortcuts
\importpaperaux{cpsv16-}{1606.00608/MPDO-22-12-17-2-tnlean}
\newcommand{\cpsvref}[1]{\paperref{cpsv16}{#1}}
\newcommand{\cpsveqref}[1]{\papereqref{cpsv16}{#1}}

% Verdict marker: \gapnote{<kind>}{<status>}. Kind is the policy
% classification (clarification, local-correction, scope-restriction,
% unfaithful, false-source, open-gap); status is open, wip (elimination
% underway), resolved, or historical. Machine-read by the site index;
% prints a verdict line.
\newcommand{\gapnote}[2]{\par\noindent\textbf{Verdict:} #1 (#2).\par}


\title{Active physical support compression in the CPSV16 commuting-form argument}
\author{}
\date{2026-07-30}

\begin{document}
\maketitle
\gapnote{scope-restriction}{open}

\section{The source passage}

In the standing Case~I setting of Cirac--P\'erez-Garc\'ia--Schuch--Verstraete,
$\mathcal K$ is an injective normal tensor that generates matrix product density
operators \cite{Cirac2016MPDO_arXiv}. Under saturation of the area law,
Lemma~\cpsvref{propSN} of Ref.~\cite{Cirac2016MPDO_arXiv} yields an isometry $U$
and factor families $l_k$ and $r_k$ such that
\begin{align}
    U\mathcal K_{\beta,\alpha}U^*
        &=
    \bigoplus_k (l_k)_\beta\otimes(r_k)_\alpha.
    \label{eq:cpsv16_active_physical_support_compression_sector_factorization}
\end{align}
It defines
\begin{align}
    \eta_{k,h}
        &=
    \sum_a (r_k)_a\otimes(l_h)_a,
    \label{eq:cpsv16_active_physical_support_compression_neighboring_operator}\\
    \eta_{k,h}
        &\geq 0,
    \label{eq:cpsv16_active_physical_support_compression_neighboring_positivity}\\
    T_{k,h}
        &=
    \tr(\eta_{k,h}).
    \label{eq:cpsv16_active_physical_support_compression_transition_matrix}
\end{align}
The source shows that the nonnegative matrix $T$ is primitive
\cite{Cirac2016MPDO_arXiv}.\footnote{Source:
\path{Papers/1606.00608/MPDO-22-12-17-2.tex}:1374--1471. The definitions of
$U$, $r_k$, $l_k$, and $\eta_{k,h}$ occur at lines~1436--1450. The definition
of $T_{k,h}$ occurs at lines~1452--1455, and its primitivity is proved through
line~1471.}

The cited passage does not construct the active physical support compression:
it chooses no isometries onto the joint column supports of the left and right
factor families, assigns no coordinate spaces to inactive sectors, and does
not assemble the compressed injective tensor.

Proposition~\cpsvref{3to4} of Ref.~\cite{Cirac2016MPDO_arXiv} uses the
uncompressed neighboring operators on the full factor spaces to construct
commuting positive bonds.\footnote{Source:
\path{Papers/1606.00608/MPDO-22-12-17-2.tex}:1570--1594.}
The subsequent Proposition~\cpsvref{prop3to4} of
Ref.~\cite{Cirac2016MPDO_arXiv} treats a simple tensor with a basis of normal
tensors, invoking the preceding injective result.\footnote{Source:
\path{Papers/1606.00608/MPDO-22-12-17-2.tex}:1786--1797.}
These two statements remain distinct in the source.

\section{The compression construction}

The formalization introduces an auxiliary compression after
Lemma~\cpsvref{propSN} and before the proof corresponding to
Proposition~\cpsvref{prop3to4} of Ref.~\cite{Cirac2016MPDO_arXiv}. The
construction is developed across five supporting modules.\footnote{The supporting
formal modules are
\doclink{TNLean/MPS/MPDO/PhysicalSectorActiveFactorSupport},
\doclink{TNLean/MPS/MPDO/PhysicalSectorActiveSupportData},
\doclink{TNLean/MPS/MPDO/PhysicalSectorActiveRestriction},
\doclink{TNLean/MPS/MPDO/PhysicalSectorActiveNeighboring}, and
\doclink{TNLean/MPS/MPDO/PhysicalSectorCoordinateTransport}.}

\begin{enumerate}
    \item For each active sector $k$, choose isometries $V^L_k$ and $V^R_k$
    whose ranges are the joint column supports of $\{l_{k,\beta}\}_\beta$ and
    $\{r_{k,\alpha}\}_\alpha$, and assign zero-dimensional coordinate spaces
    to each inactive sector. The resulting product range is the family-support
    projection of the sector product family. Injectivity of the original
    tensor and positivity of each neighboring operator together give the
    adjoint-closure identity needed for two-sided absorption by these support
    projections.

    \item Define the compressed factors by
    \begin{align}
        \widetilde l_{k,\beta}
            &=
        (V^L_k)^*l_{k,\beta}V^L_k,
        \label{eq:cpsv16_active_physical_support_compression_left_factor}\\
        \widetilde r_{k,\alpha}
            &=
        (V^R_k)^*r_{k,\alpha}V^R_k.
        \label{eq:cpsv16_active_physical_support_compression_right_factor}
    \end{align}

    \item Assemble the sector support isometries into a block-diagonal map
    and pull it back to the original physical coordinates along $U^*$, yielding
    \begin{align}
        V
            &:
        \C^e\longrightarrow\C^d.
        \label{eq:cpsv16_active_physical_support_compression_physical_inclusion}
    \end{align}
    The projection $VV^*$ is the canonical joint column support of the
    original physical slices.

    \item Define the active restriction by
    \begin{align}
        K_{\mathrm{act}}
            &=
        V^*KV.
        \label{eq:cpsv16_active_physical_support_compression_restricted_tensor}
    \end{align}
    This tensor acts on $\C^e$, where $e$ is the total dimension of the retained
    supports. Injectivity of $K$ and positivity of every neighboring operator
    ensure that $K_{\mathrm{act}}$ is injective.

    \item The compressed neighboring operator is the congruence of the
    ambient operator by the Kronecker product of the source right-support
    isometry and target left-support isometry; this congruence preserves
    positive semidefiniteness.

    \item Suppose the ambient operators satisfy the normalization of CPSV16
    Theorem~\cpsvref{thm:main-simple}(iv) \cite{Cirac2017MPDO_AnnPhys}:
    \begin{align}
        \tr(\eta_{k,h})
            &=
        a_kb_h,
        \label{eq:cpsv16_active_physical_support_compression_trace_factorization}\\
        \sum_k a_kb_k
            &=
        1.
        \label{eq:cpsv16_active_physical_support_compression_trace_normalization}
    \end{align}
    Compression preserves $\tr(\eta_{k,h})$ whenever $k$ and $h$ are active.
    For an inactive sector $k$,
    \begin{align}
        \eta_{k,k}
            &= 0,
        \label{eq:cpsv16_active_physical_support_compression_inactive_neighbor}\\
        a_kb_k
            &= 0,
        \label{eq:cpsv16_active_physical_support_compression_inactive_weight}
    \end{align}
    so restricting $a_k$ and $b_k$ to active sectors preserves
    Eqs.~\eqref{eq:cpsv16_active_physical_support_compression_trace_factorization}
    and~\eqref{eq:cpsv16_active_physical_support_compression_trace_normalization}.
    This transport is purely auxiliary and does not establish condition~(iv)
    itself.

    \item The coordinate map
    \begin{align}
        M
            &\longmapsto
        VMV^*
        \label{eq:cpsv16_active_physical_support_compression_coordinate_transport}
    \end{align}
    transports an MPO tensor together with its two-site and four-site
    closures.
\end{enumerate}

None of these compression steps appears in lines~1406--1471 of
Ref.~\cite{Cirac2016MPDO_arXiv}. While the source supplies the sector
factorization and neighboring operators, the formal development adds
compression to remove inactive dimensions before applying the theorem
corresponding to Proposition~\cpsvref{prop3to4}. Injectivity and
neighboring-operator positivity ensure that the joint column-support
projections absorb every physical slice on both sides; compression by their
adjoints therefore preserves exact two-sided recovery.

\section{Affected declarations}

The ten load-bearing declarations are:
\begin{enumerate}
    \item the support isometries for an active sector,\footnote{
    \leanid{MPOTensor.PhysicalSectorFactorization.exists_activeSector_factorSupportIsometries};}

    \item the simultaneous factor-support coordinates,\footnote{
    \leanid{MPOTensor.PhysicalSectorFactorization.ActiveFactorSupportData};}

    \item the compressed left factor,\footnote{
    \leanid{MPOTensor.PhysicalSectorFactorization.compressedLeftTensor};}

    \item recovery of an original physical slice from sector coordinates,\footnote{
    \leanid{MPOTensor.PhysicalSectorFactorization.physicalSlice_eq_conjTranspose_mul_sectorCoordinateTensor_mul};}

    \item the restricted tensor,\footnote{
    \leanid{MPOTensor.PhysicalSectorFactorization.activePhysicalSupportRestriction};}

    \item compression of the sectorwise factorization by the joint support
    inclusion,\footnote{
    \leanid{MPOTensor.PhysicalSectorFactorization.dependentSupportInclusion_compression};}

    \item exact recovery in the original physical coordinates,\footnote{
    \leanid{MPOTensor.PhysicalSectorFactorization.changePhysicalBasis_physicalSupportInclusion_activePhysicalSupportRestriction};}

    \item injectivity of the restricted tensor,\footnote{
    \leanid{MPOTensor.PhysicalSectorFactorization.activePhysicalSupportRestriction_isInjective};}

    \item the assembled restriction datum,\footnote{
    \leanid{MPOTensor.PhysicalSectorFactorization.activePhysicalSupportRestrictionData};}

    \item transport of the neighboring trace factorization.\footnote{
    \leanid{MPOTensor.PhysicalSectorFactorization.activeRestrictedPhysicalSectorFactorization_neighboringTraceFactorization}.}
\end{enumerate}

\section{Blueprint dependency}

The dependency chain is
\begin{align}
    \text{Active restriction}
        &\longrightarrow
    \text{Active neighboring}
        \longrightarrow
    \text{Active bond}
        \longrightarrow
    \text{BNT GSNNCH}.
    \label{eq:cpsv16_active_physical_support_compression_blueprint_dependency}
\end{align}
\footnote{The formal dependency path connects
\doclink{TNLean/MPS/MPDO/PhysicalSectorActiveRestriction},
\doclink{TNLean/MPS/MPDO/PhysicalSectorActiveNeighboring},
\doclink{TNLean/MPS/MPDO/PhysicalSectorActiveBond}, and
\doclink{TNLean/MPS/MPDO/GSNNCHSectorSum}.}
This chain feeds the capstone theorem for the BNT physical-sector GSNNCH
form.\footnote{Capstone declaration:
\leanid{MPOTensor.hasGSNNCHForm_of_bntLayerOrthogonal_of_physicalSectorFactorization},
labelled \path{thm:mpdo_bnt_physical_sector_gsnnch} in
\path{blueprint/src/chapter/ch21_mpdo_rfp_simple_local_structure_capstone.tex}.}
Although the compression is mathematically sound, the citation record must
reflect that this auxiliary construction is absent from the cited source
propositions.

\section{Conclusion}

The verdict is a scope restriction. The formal development first restricts to
the canonical active physical support and handles coordinate changes
separately, whereas the source propositions in Ref.~\cite{Cirac2016MPDO_arXiv}
omit these constructions. The compression itself is verified; this note records
the distinction between the source presentation and the formal proof route.

\bibliographystyle{alpha}
\bibliography{references}

\end{document}
